\documentclass[a4paper,11pt]{article}

\usepackage[T1]{fontenc}
\usepackage[utf8]{inputenc}
\usepackage{amsmath,amssymb,bm}
\usepackage[english]{babel}
\usepackage{geometry}
\geometry{a4paper,margin=2.3cm}
\usepackage{authblk}
\renewcommand*{\Affilfont}{\itshape}
\renewcommand*{\Authfont}{\bfseries}
\date{}

\begin{document}

\title{Robust nonlinear projection-based reduced-order models: comparative assessment of closure and manifold strategies}

\author[1]{S. Ares de Parga}
\author[2]{Y. Maday}

\affil[1]{Centre Internacional de M\`{e}todes Num\`{e}rics en Enginyeria (CIMNE), Barcelona, Spain}
\affil[2]{Laboratoire Jacques-Louis Lions (LJLL), Sorbonne Universit\'e and Universit\'e Paris Cit\'e, CNRS, Paris, France}

\maketitle
\thispagestyle{empty}
\vspace{-0.7cm}

\paragraph{Abstract.}
This talk presents a comparative study of ANN-enhanced projection-based reduced-order models for nonlinear parametric conservation laws.  The focus is on robust online prediction under aggressive compression, where the reduced model must remain accurate after hyper-reduction and outside the small set of high-dimensional training simulations.  We compare intrusive closure-based PROM--ANN models, an intrusive POD-autoencoder latent-manifold PROM, and two non-intrusive surrogates mapping parameter--time inputs either directly to POD coefficients or through a low-dimensional latent representation.

A central point of the study is consistency between data generation, training, and online evaluation.  All learned models are trained from coefficient trajectories generated by a linear hyper-reduced LSPG ROM in a fixed reduced space, rather than from direct orthogonal projection of high-dimensional snapshots.  The intrusive online models use the same discrete residual, Gauss--Newton LSPG solver, and ECSW hyper-reduction framework.  This builds on closure-enhanced projection ideas such as PROM--ANN \cite{barnett2023neural}, nonlinear-manifold ROM formulations \cite{aghili2026nonlinear}, hyper-reduction with ECSW \cite{grimberg2021mesh}, latent POD-DL-ROM concepts \cite{fresca2022pod}, and recent machine-learning-assisted PROM training and closure studies \cite{sibuet2025discrete,aresdeparga2026regression}.  We also use an $M$-regularized LSPG-sensitive POD metric to reduce the feedback of high-coordinate closure errors into the resolved online coordinates.

The numerical results are shown on a $250\times250$ two-dimensional parametric inviscid Burgers benchmark.  Starting from a baseline set of nine HDM trajectories, we examine the effect of enriching the learning dataset with additional low-cost linear-HPROM trajectories, including points in an expanded parameter region.  The results highlight when neural enrichment improves extrapolatory robustness, when intrusive feedback remains the limiting factor, and how closure-based and latent-manifold ROMs compare with purely non-intrusive predictors in terms of accuracy, online cost, and achieved speedup.

\begin{thebibliography}{9}
\bibitem{barnett2023neural}
J.~L. Barnett, C.~Farhat, and Y.~Maday, \emph{J. Comput. Phys.}, 492:112420, 2023.

\bibitem{aghili2026nonlinear}
J.~Aghili, H.~Ballout, Y.~Maday, and C.~Prud'Homme, \emph{arXiv preprint}, arXiv:2601.13712, 2026.

\bibitem{fresca2022pod}
S.~Fresca and A.~Manzoni, \emph{Comput. Methods Appl. Mech. Eng.}, 388:114181, 2022.

\bibitem{grimberg2021mesh}
S.~Grimberg, C.~Farhat, R.~Tezaur, and C.~Bou-Mosleh, \emph{Int. J. Numer. Meth. Eng.}, 122:1846--1874, 2021.

\bibitem{sibuet2025discrete}
N.~Sibuet, S.~Ares de~Parga, J.~R. Bravo, and R.~Rossi, \emph{Axioms}, 14(5):385, 2025.

\bibitem{aresdeparga2026regression}
S.~Ares de~Parga, R.~Tezaur, C.~G. Hern\'{a}ndez, and C.~Farhat, \emph{Comput. Methods Appl. Mech. Eng.}, 2026.
\end{thebibliography}

\end{document}
